\clearpage
\section{Gaussian states on the lattice}\label{sec:gaussianstates}

\noindent\motiv{\reve{The automata of Chapter~\ref{sec:ccrautomata} act on observables, and a
dynamics of observables becomes a dynamics of physical situations only once it acts on states. This
chapter introduces the states of the lattice CCR algebra through their characteristic functions,
singles out the Gaussian states, whose characteristic functions are exponentials of quadratic
forms, and derives the condition that the quadratic form of a Gaussian state must satisfy. Because
a Gaussian state is determined by a covariance form on the whole infinite lattice, whereas the
classical theory of Gaussian states concerns covariance matrices of finitely many modes, we devote a large
part of the chapter to showing that the two descriptions agree: restricted to a finite
region, a covariance form is a covariance matrix of finitely many modes, and the condition on the
form becomes the familiar matrix inequality on every region. On the chain the condition can also be
read off pointwise from a Fourier transform, which leads to a matrix condition that saturates the state condition at every wave number, and to an
explicit law for the evolution of the covariance under a Gaussian automaton.}}

\medskip
\noindent\motiv{\reve{As in Chapter~\ref{sec:ccrautomata}, $G$ is an abelian group, the cell carries
$n$ modes with coordinates ordered positions first, $u(x)=(q_1,\dots,q_n,p_1,\dots,p_n)$, and
$\sigma(u,w)=\sum_xu(x)^{\mathsf T}\Omega\,w(x)$ with
$\Omega=\begin{psmallmatrix}0&I_n\\-I_n&0\end{psmallmatrix}$, so that
$\Omega^{\mathsf T}=-\Omega=\Omega^{-1}$. For a site $x$ and a cell index $i$ we write $e_{x,i}\in V$
for the vector whose value at $x$ is the $i$-th unit vector and which vanishes elsewhere, so that
$\sigma(e_{x,i},e_{y,j})=\delta_{xy}\Omega_{ij}$. Since products and adjoints of Weyl operators are
multiples of Weyl operators, the linear span of the Weyl operators is a $*$-subalgebra of
$\CCR(V,\sigma)$, and it is dense by the definition of $\CCR(V,\sigma)$ as its closure. For a
complex vector or matrix, ${}^\dagger$ denotes the conjugate transpose, and $A\ge0$ means that the
Hermitian matrix $A$ is positive semidefinite.}}

\subsection{States and characteristic functions}

\noindent\motiv{\reve{A state assigns to each observable its expectation value, and the only
requirements on such an assignment are linearity, normalization and positivity on the squares
$a^*a$, which are the observables with non-negative spectrum.}}

\begin{definition}[State]\label{def:gs-state}
\reve{A \emph{state} on $\CCR(V,\sigma)$ is a continuous $\C$-linear functional $\omega$ with
$\omega(1)=1$ and $\omega(a^*a)\ge0$ for every $a$, where, for a complex number $z$, $z\ge0$ means that
$z$ is real and non-negative.}
\end{definition}

\noindent\reve{We require positivity of the complex numbers $\omega(a^*a)$, so that no order on
the algebra itself needs to be introduced. Continuity follows from positivity on every C*-algebra,
but it is included in the definition because it is the property the arguments below use directly.}

\noindent\motiv{\reve{Since the Weyl operators span a dense subalgebra, a continuous functional is
known once its values on them are known, and these values form a function on the phase space. The
products $W(u_j)^*W(u_k)$ are again multiples of Weyl operators, and positivity of the state on
combinations of such products becomes a positivity property of this function.}}

\begin{definition}[Characteristic function]\label{def:gs-chi}
\reve{The \emph{characteristic function} of a state $\omega$ is $\chi_\omega(u)=\omega(W(u))$. For a
function $\chi\colon V\to\C$ and a family $(u_j)_{j\in I}$ in $V$ indexed by a finite set $I$, the
\emph{twisted kernel} is the matrix
\[K_{jk}=e^{\ii\sigma(u_j,u_k)/2}\,\chi(u_k-u_j)\qquad(j,k\in I).\]}
\end{definition}

\begin{proposition}[Twisted positivity]\label{prop:gs-twisted}
\reve{Let $\omega$ be a state with characteristic function $\chi$.
\begin{enumerate}[label=(\alph*)]
\item For every finite family $(u_j)_{j\in I}$ and all $c\in\C^I$,
$\sum_{j,k}\bar c_jc_k\,e^{\ii\sigma(u_j,u_k)/2}\chi(u_k-u_j)\ge0$, so that the twisted kernel of
$\chi$ on every finite family is positive semidefinite.
\item $\chi(0)=1$, $\chi(-u)=\overline{\chi(u)}$ and $|\chi(u)|\le1$ for every $u\in V$.
\end{enumerate}}
\end{proposition}

\begin{proof}\leavevmode
\reve{\begin{enumerate}[label=(\alph*)]
\item By the Weyl relations,
$W(u_j)^*W(u_k)=W(-u_j)W(u_k)=e^{\ii\sigma(u_j,u_k)/2}W(u_k-u_j)$, so that for
$X=\sum_kc_kW(u_k)$ the product $X^*X$ equals
$\sum_{j,k}\bar c_jc_k\,e^{\ii\sigma(u_j,u_k)/2}W(u_k-u_j)$, and the sum in question is
$\omega(X^*X)\ge0$. A complex matrix $K$ whose quadratic form $c^\dagger Kc$ is real for every $c$
is Hermitian, because the quadratic form of $K-K^\dagger$, which is
$c^\dagger Kc-\overline{c^\dagger Kc}$, then vanishes identically, and a complex matrix with
vanishing quadratic form is zero by polarization. The twisted kernel is therefore Hermitian with
non-negative quadratic form, that is, positive semidefinite.
\item First, $\chi(0)=\omega(1)=1$. For the family $(0,u)$ the phases are trivial, because
$\sigma(0,u)=\sigma(u,0)=0$, and the twisted kernel is
$\begin{psmallmatrix}1&\chi(u)\\\chi(-u)&1\end{psmallmatrix}$. By (a) it is Hermitian, which gives
$\chi(-u)=\overline{\chi(u)}$, and its determinant $1-|\chi(u)|^2$ is non-negative.
\end{enumerate}}
\end{proof}

\begin{remark}[Twisted positive-definiteness]\label{rem:gs-bochner}
\reve{Part~(a) says that $\chi$ is \emph{$\sigma$-positive-definite}. For $\sigma=0$ this is the
positive-definiteness of Bochner's theorem, by which a continuous positive-definite function on a
finite-dimensional real vector space with $\chi(0)=1$ is the characteristic function of a
probability measure, and the twisted condition is its analogue for the canonical commutation
relations. The analogue of Bochner's theorem itself, for finitely many degrees of freedom, is the
quantum Bochner theorem of Kastler~\cite{kastler} and of Loupias and Miracle-Sole~\cite{lms}: a
continuous $\sigma$-positive-definite function on a finite-dimensional symplectic space with
$\chi(0)=1$ is the characteristic function of a density operator in the Schr\"odinger
representation, and its symplectic Fourier transform is the Wigner function of that density
operator, just as the Fourier transform of $\chi$ is the probability density in the commutative
case} \vetsource{(statement recalled, not checked against the two sources)}\reve{. Both hypotheses
of that theorem belong to finitely many degrees of freedom, continuity of $\chi$ and a Schr\"odinger
representation in which the density operator lives. On the infinite lattice there is no such
distinguished representation, and a state need not be given by a density operator in any fixed
representation, as Remark~\ref{rem:gs-existence} shows for the representation on $\ell^2(V)$. The
result that replaces the quantum Bochner theorem there is the characterization of states by their
characteristic functions quoted in that remark, which requires no continuity of $\chi$.}
\revs{This Fourier transform acts on the phase-space variable $u$, unlike the lattice Fourier
transform $x\mapsto k$ of Subsection~\ref{sec:gs-fourier}, which acts on the site variable of a
vector.}
\end{remark}

\noindent\motiv{\reve{Two properties of characteristic functions suffice for the evolution of
states: a state is determined by its characteristic function, and pulling a state back along an
automorphism that relabels the Weyl operators composes the characteristic function with the
relabelling.}}

\begin{proposition}[Pullback along automorphisms]\label{prop:gs-pullback}
\reve{\begin{enumerate}[label=(\alph*)]
\item Two states with the same characteristic function are equal.
\item If $\alpha$ is a $*$-automorphism of $\CCR(V,\sigma)$ and $\omega$ a state, then
$\omega\circ\alpha$ is a state. If in addition $\alpha(W(u))=W(Qu)$ for a map $Q\colon V\to V$, then
$\chi_{\omega\circ\alpha}=\chi_\omega\circ Q$.
\end{enumerate}}
\end{proposition}

\begin{proof}\leavevmode
\reve{\begin{enumerate}[label=(\alph*)]
\item Two continuous linear functionals that agree on the Weyl operators agree on their linear
span, and hence on its closure, which is the whole algebra.
\item A $*$-automorphism of a C*-algebra is isometric, so that $\omega\circ\alpha$ is continuous
as well as linear; it is normalized because $\alpha(1)=1$, and positive because
$\omega(\alpha(a^*a))=\omega(\alpha(a)^*\alpha(a))\ge0$. Finally
$\chi_{\omega\circ\alpha}(u)=\omega(W(Qu))=\chi_\omega(Qu)$.
\end{enumerate}}
\end{proof}

\begin{corollary}[Translation invariance]\label{cor:gs-transl}
\reve{A state $\omega$ satisfies $\omega\circ\alpha_a=\omega$ for every $a\in G$ if and only if
$\chi_\omega\circ\kappa_{a*}=\chi_\omega$ for every $a\in G$.}
\end{corollary}

\begin{proof}
\reve{By Proposition~\ref{prop:gs-pullback}(b) with $Q=\kappa_{a*}$, the state $\omega\circ\alpha_a$
has the characteristic function $\chi_\omega\circ\kappa_{a*}$, and by
Proposition~\ref{prop:gs-pullback}(a) it equals $\omega$ exactly when the two characteristic
functions agree.}
\end{proof}

\subsection{Gaussian states}

\noindent\motiv{\reve{In classical probability a Gaussian measure is one whose characteristic
function is the exponential of a quadratic polynomial, and the quantum definition takes over this
description word for word; on the CCR algebra these are the quasi-free states of Manuceau and
Verbeure~\cite{manuceau-verbeure}. The factor $\tfrac14$ in the exponent is the normalization in which the
vacuum of a single oscillator has covariance $1$, as Example~\ref{ex:gs-vacuum} below shows; it
agrees with the finite-mode convention in which the covariance matrix is twice the matrix of
symmetrized second moments.}}

\begin{definition}[Gaussian state]\label{def:gs-gaussian}
\reve{A \emph{Gaussian state} is a state $\omega$ together with a symmetric bilinear form $\gamma$ on
$V$, its \emph{covariance form}, and a linear form $d$ on $V$, its \emph{displacement}, such that
\[\chi_\omega(u)=\exp\bigl(\ii\,d(u)-\tfrac14\gamma(u,u)\bigr)\qquad\text{for all }u\in V.\]}
\end{definition}

\noindent\reve{Every result below concerns a Gaussian state given together with its data
$(\gamma,d)$. That a state with prescribed data exists is not proved in these notes, and no result
below depends on it; the examples of the chapter are accordingly covariance forms satisfying the
state condition of Definition~\ref{def:gs-statecond}, which by Theorem~\ref{thm:gs-necessity} is
necessary for a Gaussian state, and Remark~\ref{rem:gs-existence} explains why existence cannot
simply be read off from the construction of the algebra.}

\begin{remark}[The data are determined by the state]\label{rem:gs-data}
\reve{Although the definition lists $\gamma$ and $d$ as data, a state admits at most one pair of
them. Indeed, $|\chi_\omega(u)|=\exp(-\tfrac14\gamma(u,u))$ determines $\gamma(u,u)$ for every $u$,
and hence the symmetric form $\gamma$ by polarization; then
$\chi_\omega(tu)/|\chi_\omega(tu)|=e^{\ii t\,d(u)}$ for all $t\in\R$, and two real numbers $r,r'$
with $e^{\ii tr}=e^{\ii tr'}$ for every real $t$ are equal, as one sees by taking
$|t(r-r')|<2\pi$ with $t\neq0$, so that $d(u)$ is determined as well.}
\end{remark}

\begin{remark}[Normalizations]\label{rem:gs-normalization}
\reve{In the normalization of these notes $\chi_\omega(u)=\exp(\ii d(u)-\tfrac14\gamma(u,u))$, and
for one mode the corresponding covariance matrix $\gamma_B$ is twice the matrix of symmetrized
second moments, as the paragraph before Example~\ref{ex:gs-squeezed} shows. Other normalizations of the literature write the exponent with a
different constant, for instance as $-\tfrac12\gamma'(u,u)$, and then $\gamma=2\gamma'$; such a
rescaling of $\gamma$ by a constant makes conditions such as $\sigma(u,w)^2\le\gamma(u,u)\gamma(w,w)$
acquire factors, here $\tfrac14\sigma(u,w)^2\le\gamma'(u,u)\gamma'(w,w)$, without changing their
content. The sign of $\ii\Omega$ is likewise a matter of convention: for a real symmetric matrix
$\gamma_B$ of $M$ modes the matrices $\gamma_B-\ii\Omega_M$ and $\gamma_B+\ii\Omega_M$, with
$\Omega_M$ the symplectic matrix of \revt{Subsection}~\ref{sec:gs-regions}, are complex conjugates of each
other, and the complex conjugate of a positive semidefinite matrix is positive semidefinite, so
that $\gamma_B-\ii\Omega_M\ge0$ and $\gamma_B+\ii\Omega_M\ge0$ are equivalent and either may be called
the finite-mode condition.}
\end{remark}

\begin{remark}[Existence and regularity]\label{rem:gs-existence}
\reve{The representation of $\CCR(V,\sigma)$ on $\ell^2(V)$ by which the algebra was constructed
contains no Gaussian state. For $u\neq0$ and $t\neq0$, the vector $W(tu)\delta_v$ is a multiple of
$\delta_{v-tu}$, which is orthogonal to $\delta_v$, so that $\|W(tu)\delta_v-\delta_v\|=\sqrt2$ and
$t\mapsto W(tu)$ is not strongly continuous, and the representation is therefore not regular.
Let $\rho=\sum_ip_i|\psi_i\rangle\langle\psi_i|$ be a density operator on $\ell^2(V)$ and $u\neq0$. Each
$\psi_i$ has countable support $\operatorname{supp}\psi_i$, and
$\langle\psi_i,W(tu)\psi_i\rangle=\sum_v\overline{\psi_i(v)}\,e^{\ii\sigma(tu,v)/2}\psi_i(v+tu)$
vanishes unless $tu\in\operatorname{supp}\psi_i-\operatorname{supp}\psi_i$, which happens for only countably many $t$. Hence
$\operatorname{tr}(\rho\,W(tu))=0$ for all but countably many $t\in\R$, whereas a Gaussian
characteristic function vanishes nowhere, and no normal state of this representation is Gaussian.}

\reve{That a Gaussian function with a covariance form satisfying the state condition nevertheless
defines a state follows from the theorem that a function $\chi$ on $V$ with $\chi(0)=1$ and positive
semidefinite twisted kernels is the characteristic function of a unique state on the CCR
algebra~\cite{mstv},}~\cite[\S5.2]{bratteli-robinson}\reve{. Its proof passes through
the GNS construction on the group algebra of $V$ twisted by $\sigma$ and through the uniqueness of
the CCR algebra, and we do not reproduce it here. The uniqueness of the CCR algebra, and with it the
theorem, requires the form $\sigma$ to be nondegenerate, and this is the one place where
nondegeneracy enters the existence question: if $\sigma(u,w)=0$ for every $w\in V$, then pairing
with $w=e_{x,i}$ gives $u(x)^{\mathsf T}\Omega=0$, hence $u(x)=0$, at every site, so that $u=0$.
The elementary half, that the state condition makes
the Gaussian twisted kernels positive semidefinite, is Remark~\ref{rem:gs-converse}.}
\end{remark}

\subsection{The state condition and its necessity}

\noindent\motiv{\reve{For finitely many modes, a real symmetric matrix $\gamma_B$ is the covariance matrix of a
Gaussian state exactly when $\gamma_B-\ii\Omega_M\ge0$ (equivalently $\gamma_B+\ii\Omega_M\ge0$), a form of the uncertainty
principle. On the lattice there is no matrix, but the condition can be imposed on every finite
family of vectors, and on a real vector space it can be rewritten without complex numbers in two
ways, the second of which makes it evident that the sign in front of $\ii\sigma$ plays no
role.}}

\begin{definition}[State condition]\label{def:gs-statecond}
\reve{Let $\gamma$ be a symmetric and $\sigma$ an alternating bilinear form on a real vector space
$V$. The pair $(\gamma,\sigma)$ satisfies
\begin{enumerate}[label=(\alph*)]
\item the \emph{state condition} if for every finite family $(u_j)_{j\in I}$ in $V$ the complex
matrix $\bigl(\gamma(u_j,u_k)+\ii\sigma(u_j,u_k)\bigr)_{j,k}$ is positive semidefinite;
\item the state condition \emph{in pair form} if $\gamma(a,a)+\gamma(b,b)\ge2\sigma(a,b)$ for all $a,b\in V$;
\item the state condition \emph{in real-pair form} if $\gamma(u,u)\ge0$ and $\sigma(u,w)^2\le\gamma(u,u)\,\gamma(w,w)$
for all $u,w\in V$.
\end{enumerate}}
\end{definition}

\noindent\reve{Condition~(a) can equally be read on complex vectors: if $\gamma_\C$ and $\sigma_\C$
denote the sesquilinear extensions to $V_\C=V\otimes\C$, then for $\mu=\sum_kc_ku_k$ the quadratic
form of the matrix in~(a) at $c$ is $\gamma_\C(\mu,\mu)+\ii\sigma_\C(\mu,\mu)$, and every
$\mu=a+\ii b$ arises in this way from the family $(a,b)$ with $c=(1,\ii)$.}

\begin{proposition}[Equivalent forms]\label{prop:gs-statecond}
\reve{For symmetric $\gamma$ and alternating $\sigma$, the three conditions of
Definition~\ref{def:gs-statecond} are equivalent, and they are also equivalent to condition~(a)
with $\gamma-\ii\sigma$ in place of $\gamma+\ii\sigma$.}
\end{proposition}

\begin{proof}
\reve{The matrix in (a) is Hermitian, since $\gamma$ is symmetric and $\sigma$ antisymmetric, and we
first compute its quadratic form. For the family $(a,b)$ and $c=(1,\ii)$, the diagonal contributes
$\gamma(a,a)+\gamma(b,b)$, since $\sigma(a,a)=\sigma(b,b)=0$, while the off-diagonal terms
contribute $\ii\bigl(\gamma(a,b)-\gamma(b,a)\bigr)=0$ from $\gamma$ and
$\ii\cdot\ii\bigl(\sigma(a,b)-\sigma(b,a)\bigr)=-2\sigma(a,b)$ from $\sigma$. For a general family
and general coefficients, bilinearity reduces the computation to this case with
$a=\sum_k(\operatorname{Re}c_k)u_k$ and $b=\sum_k(\operatorname{Im}c_k)u_k$, and therefore
\begin{equation}\label{eq:gs-pair}
\sum_{j,k}\bar c_jc_k\,(\gamma+\ii\sigma)(u_j,u_k)=\gamma(a,a)+\gamma(b,b)-2\sigma(a,b).
\end{equation}
The quadratic forms in (a) are thus non-negative for all families and coefficients exactly when
(b) holds, which proves (a)$\Leftrightarrow$(b).}

\reve{Assume (b). Taking $b=0$ gives $\gamma(a,a)\ge0$. Given $u,w$ and $s>0$, (b) applied to $a=su$
and $b=\pm w/s$ gives $s^2\gamma(u,u)+s^{-2}\gamma(w,w)\ge2|\sigma(u,w)|$. When $\gamma(u,u)$ and
$\gamma(w,w)$ are both positive, the choice $s^2=\sqrt{\gamma(w,w)/\gamma(u,u)}$ yields
$\sqrt{\gamma(u,u)\gamma(w,w)}\ge|\sigma(u,w)|$, which is (c); when $\gamma(u,u)=0$, the inequality
$s^{-2}\gamma(w,w)\ge2|\sigma(u,w)|$ holds for every $s>0$ and forces $\sigma(u,w)=0$ as
$s\to\infty$, so that (c) holds in the form $0\le0$, and the case $\gamma(w,w)=0$ is the same with
$s\to0$. Conversely, (c) and the inequality between the geometric and the arithmetic mean give
$2\sigma(a,b)\le2\sqrt{\gamma(a,a)\gamma(b,b)}\le\gamma(a,a)+\gamma(b,b)$, which is (b).}

\reve{Condition (c) does not change when $\sigma$ is replaced by $-\sigma$, and condition (a) for
$-\sigma$ is condition (a) with $\gamma-\ii\sigma$, which proves the last statement.}
\end{proof}

\noindent\motiv{\reve{The necessity of the state condition is the statement that positivity of a
Gaussian state forces the uncertainty relation on its covariance form. The argument below uses
positivity only on elements of the form \revt{$\sum_kb_k\bigl(\delta_k^{-1}W(su_k)-1\bigr)$} and lets $s$ tend to zero; the
subtraction of the identity removes the constant term of the exponential, which would otherwise
dominate, and the limit is controlled by an explicit bound, so that no derivatives are needed.}}

\begin{theorem}[Necessity of the state condition]\label{thm:gs-necessity}
\reve{The covariance form $\gamma$ of every Gaussian state satisfies the state condition with
$\sigma$, in each of the three equivalent forms of Definition~\ref{def:gs-statecond}.}
\end{theorem}

\begin{proof}
\reve{Let $(\omega,\gamma,d)$ be a Gaussian state, $u_1,\dots,u_J\in V$ and $b\in\C^J$, and put
$H_{jk}=\tfrac12\bigl(\gamma(u_j,u_k)+\ii\sigma(u_j,u_k)\bigr)$. The matrix $H$ is Hermitian, so that
$h=\sum_{j,k}\bar b_jb_kH_{jk}$ is real, and we have to show that $h\ge0$.}

\reve{For $s>0$ consider the family $v_0=0$ and $v_k=su_k$ for $1\le k\le J$, and put $\delta_0=1$
and $\delta_k=\exp\bigl(\ii s\,d(u_k)-\tfrac14s^2\gamma(u_k,u_k)\bigr)$. Expanding
$\gamma(v_k-v_j,v_k-v_j)=\gamma(v_j,v_j)+\gamma(v_k,v_k)-2\gamma(v_j,v_k)$ by symmetry and using the
linearity of $d$, the twisted kernel of $\chi_\omega$ on this family factorizes as
\[K_{jk}=e^{\ii\sigma(v_j,v_k)/2}\exp\bigl(\ii d(v_k-v_j)-\tfrac14\gamma(v_k-v_j,v_k-v_j)\bigr)
=\overline{\delta_j}\,\delta_k\,e^{s^2H_{jk}},\]
where $H_{0k}=H_{j0}=0$. We apply Proposition~\ref{prop:gs-twisted}(a) with the coefficients
$c_0=-B$, where $B=\sum_{k\ge1}b_k$, and $c_k=b_k/\delta_k$ for $k\ge1$, which is possible because
$\delta_k\neq0$. Then $c_k\delta_k=b_k$ for $k\ge1$ and $c_0\delta_0=-B$, and the terms of the
quadratic form with $j=k=0$, with $j=0<k$, with $k=0<j$ and with $j,k\ge1$ contribute $|B|^2$,
$-|B|^2$, $-|B|^2$ and $\sum_{j,k\ge1}\bar b_jb_ke^{s^2H_{jk}}$ respectively. Since
$\sum_{j,k\ge1}\bar b_jb_k=|B|^2$, this gives, with $\tau=s^2$,
\begin{equation}\label{eq:gs-schur}
0\le\sum_{j,k=1}^J\bar b_jb_k\bigl(e^{\tau H_{jk}}-1\bigr)\qquad\text{for every }\tau>0,
\end{equation}
which in terms of the algebra is $\omega(X^*X)\ge0$ for \revt{$X=\sum_kb_k\bigl(\delta_k^{-1}W(su_k)-1\bigr)$, the element with the coefficients $c_0,\dots,c_J$ above}.}

\reve{For a complex number $z$ with $|z|\le1$ one has
$|e^z-1-z|\le\sum_{r\ge2}|z|^r/r!\le|z|^2$. Choose $\tau_0>0$ with $\tau_0|H_{jk}|\le1$ for all
$j,k$, and let $0<\tau\le\tau_0$. Writing $e^{\tau H_{jk}}-1=\tau H_{jk}+\rho_{jk}$ with
$|\rho_{jk}|\le\tau^2|H_{jk}|^2$ and taking real parts in~\eqref{eq:gs-schur}, we obtain
\[0\le\tau h+\tau^2C,\qquad C=\sum_{j,k}|b_j|\,|b_k|\,|H_{jk}|^2,\]
so that $h\ge-\tau C$ for every $\tau\in(0,\tau_0]$, and hence $h\ge0$. Since $2h$ is the quadratic
form of condition (a) of Definition~\ref{def:gs-statecond}, this proves (a), and the other two forms
follow from Proposition~\ref{prop:gs-statecond}.}
\end{proof}

\begin{remark}[The converse for twisted kernels]\label{rem:gs-converse}
\reve{Conversely, if $\gamma$ is symmetric and satisfies the state condition, and $d$ is linear, then
the Gaussian function $\chi(u)=\exp\bigl(\ii d(u)-\tfrac14\gamma(u,u)\bigr)$ has positive
semidefinite twisted kernels. By the factorization in the proof above with $s=1$, the kernel on a
family $(u_k)$ is $K=D^\dagger ED$ with $D=\operatorname{diag}(\delta_k)$ and $E_{jk}=e^{H_{jk}}$,
and $E=\sum_{r\ge0}H^{\circ r}/r!$ is a convergent sum of entrywise powers of the positive
semidefinite matrix $H$, each of which is positive semidefinite by the Schur product theorem.
Together with the theorem cited in Remark~\ref{rem:gs-existence}, this yields a Gaussian state for
every symmetric $\gamma$ satisfying the state condition and every linear $d$.}
\end{remark}

\subsection{Finite regions and the covariance matrix}\label{sec:gs-regions}

\noindent\motiv{\reve{A finite region $O\subseteq G$ carries $n_O=n|O|$ modes, and a Gaussian state
of finitely many modes is described by a real $2n_O\times2n_O$ covariance matrix $\gamma_B$, which
is twice the matrix of symmetrized second moments of the canonical coordinates, ordered with all
positions before all momenta. We call such a matrix \emph{admissible} if it is symmetric and
$\gamma_B-\ii\Omega_{n_O}\ge0$, where
$\Omega_{M}=\begin{psmallmatrix}0&I_M\\-I_M&0\end{psmallmatrix}$ is the symplectic matrix of $M$
modes in the same order. The covariance form of a lattice state, restricted to a finite region, is
exactly such a covariance matrix, and the state condition on the region is exactly admissibility,
once we account for two differences of bookkeeping: the site coordinates of a region come site
by site, whereas the finite-mode convention puts all positions first, and the form $\gamma$ is evaluated on the labels $u$ of the Weyl operators
$W(u)$, as a two-point function is evaluated on test functions, whereas $\gamma_B$ is a matrix of
moments of the canonical coordinates, the two being related by a conjugation with $\Omega$.}}

\noindent\reve{We write $\loc(O)$ for the subspace of vectors supported in $O$, as in
\revt{the table of notation}, and say that $\gamma$ satisfies the state condition \emph{on $O$} if
the restrictions of $\gamma$ and $\sigma$ to $\loc(O)$ satisfy it; by
Proposition~\ref{prop:gs-statecond} this may be taken in the real-pair form, $\gamma(u,u)\ge0$ and
$\sigma(u,w)^2\le\gamma(u,u)\gamma(w,w)$ for $u,w\in\loc(O)$.}

\begin{definition}[Covariance matrix of a region]\label{def:gs-covmat}
\reve{Let $O\subseteq G$ be finite and fix an enumeration $r\colon O\times\{1,\dots,n\}\to
\{1,\dots,n_O\}$ of its modes. The \emph{site vectors} of $O$ in the finite-mode order are
\[f_{q_{r(x,i)}}=e_{x,i},\qquad f_{p_{r(x,i)}}=e_{x,n+i}\qquad(x\in O,\ 1\le i\le n),\]
so that the position coordinate of the mode $r(x,i)$ belongs to the $i$-th position of the site $x$,
and its momentum coordinate to the $i$-th momentum. For a bilinear form $\beta$ on $V$, the
\emph{restricted matrix} $R_O(\beta)$ is the real $2n_O\times2n_O$ matrix with entries
$\beta(f_a,f_b)$, rows and columns ordered $q_1,\dots,q_{n_O},p_1,\dots,p_{n_O}$. The
\emph{covariance matrix} of $\gamma$ on $O$ is
\[\cov_O(\gamma)=\Omega_{n_O}^{\mathsf T}\,R_O(\gamma)\,\Omega_{n_O}.\]}
\end{definition}

\noindent\reve{The matrix $R_O(\gamma)$ is the matrix of $\gamma$ on the site vectors of $O$ taken
site by site, conjugated by the permutation matrix that puts positions first, and the conjugation
by $\Omega_{n_O}$ is the passage from forms on labels to moments, $\gamma_B=\Omega^{\mathsf
T}\gamma\,\Omega$, which the paragraph before Example~\ref{ex:gs-squeezed} derives for one mode. A different enumeration
of the modes conjugates $\cov_O(\gamma)$ by $P\oplus P$ for a permutation matrix $P$ of the modes;
since $(P\oplus P)\Omega_{n_O}(P\oplus P)^{\mathsf T}=\Omega_{n_O}$, this preserves admissibility, and
none of the statements below depends on the enumeration.}

\begin{lemma}[Coordinates on a region]\label{lem:gs-reorder}
\reve{Let $O$ be finite.
\begin{enumerate}[label=(\alph*)]
\item $R_O(\sigma)=\Omega_{n_O}$, and $\Omega_{n_O}^{\mathsf T}\Omega_{n_O}=1$.
\item Every $u\in\loc(O)$ is $u=\sum_a\xi_af_a$ for the coordinate vector $\xi\in\R^{2n_O}$ with $\xi_a$ the corresponding entry $u(x)_i$,
and for $w=\sum_a\eta_af_a$ one has $\beta(u,w)=\xi^{\mathsf T}R_O(\beta)\,\eta$ for every bilinear
form $\beta$.
\item For $n=1$ and a one-point region $\{x\}$, $\cov_{\{x\}}(\gamma)=\Omega^{\mathsf
T}R_{\{x\}}(\gamma)\,\Omega$, where $R_{\{x\}}(\gamma)$ is the matrix of $\gamma$ on the site
vectors $e_{x,1},e_{x,2}$; and $\Omega^{\mathsf
T}\operatorname{diag}(\alpha,\beta)\,\Omega=\operatorname{diag}(\beta,\alpha)$.
\end{enumerate}}
\end{lemma}

\begin{proof}\leavevmode
\reve{\begin{enumerate}[label=(\alph*)]
\item Since $\sigma(e_{x,i},e_{y,j})=\delta_{xy}\Omega_{ij}$, the only nonzero values of $\sigma$ on
the site vectors are $\sigma(f_{q_r},f_{p_r})=1$ and $\sigma(f_{p_r},f_{q_r})=-1$, which are the
entries of $\Omega_{n_O}$. The second statement is the block computation
$\begin{psmallmatrix}0&-I\\I&0\end{psmallmatrix}\begin{psmallmatrix}0&I\\-I&0\end{psmallmatrix}=1$.
\item The vectors $f_a$ are the vectors $e_{x,i}$ with $x\in O$, each listed once, and these form a
basis of $\loc(O)$, so that $u$ is the sum of its entries times these vectors; the formula for $\beta$ is bilinearity.
\item For one site and one mode, $f_{q_1}=e_{x,1}$ and $f_{p_1}=e_{x,2}$, and $\Omega_1=\Omega$.
The identity is the product
$\begin{psmallmatrix}0&-1\\1&0\end{psmallmatrix}\begin{psmallmatrix}0&\alpha\\-\beta&0\end{psmallmatrix}
=\operatorname{diag}(\beta,\alpha)$.
\end{enumerate}}
\end{proof}

\begin{proposition}[Restriction to a region]\label{prop:gs-admissible}
\reve{Let $\gamma$ be a symmetric bilinear form on $V$ and $O$ a finite region. Then $\gamma$
satisfies the state condition on $O$ if and only if $\cov_O(\gamma)$ is admissible. In particular,
the covariance form of a Gaussian state has an admissible covariance matrix on every finite
region.}
\end{proposition}

\begin{proof}
\reve{By Lemma~\ref{lem:gs-reorder}(b), the coordinates identify $\loc(O)$ with $\R^{2n_O}$ and carry
$\gamma$ and $\sigma$ to the forms of the matrices $R_O(\gamma)$ and $R_O(\sigma)=\Omega_{n_O}$. A
complex combination of vectors of $\loc(O)$ has a complex coordinate vector $\zeta$, and by
\eqref{eq:gs-pair} the quadratic form of condition (a) of Definition~\ref{def:gs-statecond} at it is
$\zeta^\dagger\bigl(R_O(\gamma)+\ii\Omega_{n_O}\bigr)\zeta$. The state condition on $O$ is therefore
equivalent to $R_O(\gamma)+\ii\Omega_{n_O}\ge0$ and, by the last statement of
Proposition~\ref{prop:gs-statecond}, to $R_O(\gamma)-\ii\Omega_{n_O}\ge0$. Since $R_O(\gamma)$ is
symmetric, this says that $R_O(\gamma)$ is admissible.}

\reve{It remains to pass from $R_O(\gamma)$ to $\cov_O(\gamma)$. If $S$ is an invertible real matrix
with $S\Omega_{n_O}S^{\mathsf T}=\Omega_{n_O}$, then
$S(A-\ii\Omega_{n_O})S^{\mathsf T}=SAS^{\mathsf T}-\ii\Omega_{n_O}$, and since a congruence by an
invertible matrix preserves positive semidefiniteness in both directions, $A$ is admissible if and
only if $SAS^{\mathsf T}$ is. The matrix $S=\Omega_{n_O}^{\mathsf T}$ is of this kind, because
$\Omega_{n_O}^{\mathsf T}\Omega_{n_O}\Omega_{n_O}=\Omega_{n_O}$ by Lemma~\ref{lem:gs-reorder}(a),
and $SR_O(\gamma)S^{\mathsf T}=\cov_O(\gamma)$. The last statement follows from
Theorem~\ref{thm:gs-necessity}, since the state condition on $V$ implies it on $\loc(O)$.}
\end{proof}

\noindent\motiv{\reve{The familiar operations on covariance matrices of finitely many modes, the
passage to a subsystem and the direct sum of independent systems, correspond on the lattice to the
passage to a smaller region and to the union of two regions without correlations between them.}}

\begin{corollary}[Sub-regions and unions]\label{cor:gs-regions}
\reve{Let $\gamma$ be a symmetric bilinear form on $V$.
\begin{enumerate}[label=(\alph*)]
\item If $\gamma$ satisfies the state condition on a finite region $O$, it satisfies it on every
$O'\subseteq O$.
\item Let $O_1$ and $O_2$ be disjoint finite regions with $\gamma(u,w)=0$ for $u\in\loc(O_1)$ and
$w\in\loc(O_2)$. Then $\gamma$ satisfies the state condition on $O_1\cup O_2$ if and only if it
satisfies it on $O_1$ and on $O_2$.
\end{enumerate}}
\end{corollary}

\begin{proof}\leavevmode
\reve{\begin{enumerate}[label=(\alph*)]
\item The real-pair form on $\loc(O)$ implies it on the subspace $\loc(O')\subseteq\loc(O)$.
\item The direction ($\Rightarrow$) is (a). Conversely, every $a\in\loc(O_1\cup O_2)$ splits as
$a=a_1+a_2$ with $a_i\in\loc(O_i)$, and similarly $b=b_1+b_2$. The cross terms of $\gamma$ vanish
by hypothesis and by symmetry, and those of $\sigma$ vanish because $\sigma(u,w)=0$ whenever $u$ and
$w$ have disjoint supports. Hence
\[\gamma(a,a)+\gamma(b,b)-2\sigma(a,b)=\sum_{i=1,2}\bigl(\gamma(a_i,a_i)+\gamma(b_i,b_i)
-2\sigma(a_i,b_i)\bigr),\]
and the pair form on $O_1$ and on $O_2$ gives it on $O_1\cup O_2$; by
Proposition~\ref{prop:gs-statecond}, applied to the restrictions, the pair form is equivalent to the
state condition.
\end{enumerate}}
\end{proof}

\noindent\reve{In terms of covariance matrices, (a) says that $\cov_{O'}(\gamma)$, which is the
submatrix of $\cov_O(\gamma)$ on the modes of $O'$ when the enumerations are chosen compatibly, is
admissible whenever $\cov_O(\gamma)$ is, and (b) says that under the hypothesis $\cov_O(\gamma)$ is,
up to a permutation of modes, the direct sum of $\cov_{O_1}(\gamma)$ and $\cov_{O_2}(\gamma)$, which
is admissible exactly when both summands are.}

\begin{remark}[Definiteness]\label{rem:gs-posdef}
\reve{If $\gamma$ is symmetric and satisfies the state condition on a finite region $O$, then
$\gamma(u,u)>0$ for every nonzero $u\in\loc(O)$, so that $R_O(\gamma)$ and the congruent matrix
$\cov_O(\gamma)$ are positive definite. Indeed, if $\gamma(u,u)=0$, the real-pair form gives $\sigma(u,w)=0$ for every
$w\in\loc(O)$, and pairing with $w=e_{x,i}$ for $x\in O$ shows that $\Omega^{\mathsf T}u(x)=0$, so
that $u(x)=0$ for every $x\in O$ and $u=0$.}
\end{remark}

\noindent\motiv{\reve{The two basic examples of the finite-mode theory, the vacuum and the squeezed
states, are products over the sites. We give each by its covariance form, which by
Theorem~\ref{thm:gs-necessity} and Proposition~\ref{prop:gs-admissible} is what a Gaussian state
with these correlations must have.}}

\begin{example}[The vacuum]\label{ex:gs-vacuum}
\reve{The \emph{vacuum form} is $\gamma(u,w)=\sum_xu(x)^{\mathsf T}w(x)$. It satisfies the state
condition on $V$: since $\Omega$ is orthogonal,
\[\gamma(a,a)+\gamma(b,b)-2\sigma(a,b)=\sum_x\bigl|a(x)-\Omega\,b(x)\bigr|^2\ge0,\]
which is the pair form. On every finite region $R_O(\gamma)=1$, since the site vectors are
orthonormal, and hence $\cov_O(\gamma)=\Omega_{n_O}^{\mathsf T}\Omega_{n_O}=1$; for one mode and one
site this is the covariance matrix $\operatorname{diag}(1,1)$ of the ground state of the harmonic
oscillator in natural units, whose characteristic function is $\exp\bigl(-\tfrac14|u|^2\bigr)$.}
\end{example}

\noindent\motiv{\reve{Let $n=1$. On one mode, \revt{$W(u)=\exp\bigl(\ii(u_p\hat q-u_q\hat p)\bigr)$, with position and momentum operators $\hat q$ and $\hat p$,} realizes the Weyl relations of
these notes, and a centred Gaussian state with covariance matrix $\gamma_B$ has
$\chi(u)=\exp\bigl(-\tfrac14(\Omega u)^{\mathsf T}\gamma_B(\Omega u)\bigr)$, since $\Omega
u=(u_p,-u_q)$. Comparing with Definition~\ref{def:gs-gaussian} gives $\gamma=\Omega^{\mathsf
T}\gamma_B\Omega$, or equivalently $\gamma_B=\Omega\gamma\Omega^{\mathsf T}=\Omega^{\mathsf
T}\gamma\,\Omega$, the last because $\Omega^{\mathsf T}=-\Omega$. For the vacuum both sides are the
identity, and it is the squeezed state that makes the conjugation visible: for
$\gamma_B=\operatorname{diag}(\vartheta,1/\vartheta)$ with $\vartheta>0$ the characteristic function is
$\exp\bigl(-\tfrac14(u_q^2/\vartheta+\vartheta u_p^2)\bigr)$, so that $\gamma=\operatorname{diag}(1/\vartheta,\vartheta)$ on one
site.}}

\begin{example}[Squeezed product states]\label{ex:gs-squeezed}
\reve{Let $n=1$ and $\vartheta>0$. The \emph{squeezed form} is
\[\gamma(u,w)=\sum_xu(x)^{\mathsf T}\operatorname{diag}(1/\vartheta,\vartheta)\,w(x).\]
By Lemma~\ref{lem:gs-reorder}(c) its covariance matrix on a one-point region is
$\operatorname{diag}(\vartheta,1/\vartheta)$. It satisfies the state condition on $V$, because
for $a(x)=(q,p)$ and $b(x)=(q',p')$ the contribution of the site $x$ to the pair form is
\[\frac{q^2+q'^2}{\vartheta}+\vartheta(p^2+p'^2)-2(qp'-pq')
=\Bigl(\frac{q}{\sqrt\vartheta}-\sqrt\vartheta\,p'\Bigr)^2+\Bigl(\sqrt\vartheta\,p+\frac{q'}{\sqrt\vartheta}\Bigr)^2\ge0 .\]}
\end{example}

\subsection{The Fourier criterion}\label{sec:gs-fourier}

\noindent\motiv{\reve{From here on $G=\Z$. A translation-invariant covariance form is determined by
the correlations between the site $0$ and every other site, a family of $2n\times2n$ matrices, just
as a rule is determined by its coefficients. The Fourier transform then diagonalizes the
translations, and the state condition, which involves arbitrarily large families of vectors,
reduces to a condition on one $2n\times2n$ Hermitian matrix at each wave number~$k$.}}

\begin{definition}[Cell coefficients and symbol]\label{def:gs-cells}
\reve{A family $(\gamma_x)_{x\in\Z}$ of real $2n\times2n$ matrices is \emph{cell data} if
it is absolutely summable, $\sum_x\|\gamma_x\|<\infty$, and symmetric, $\gamma_{-x}=\gamma_x^{\mathsf
T}$. Its \emph{covariance form} and its \emph{symbol} are
\[\gamma(u,w)=\sum_{x,y\in\Z}u(x)^{\mathsf T}\gamma_{x-y}\,w(y),\qquad
\hat\gamma(k)=\sum_{x\in\Z}\gamma_x\,e^{-\ii kx}\quad(k\in\R).\]
The family has \emph{range} $\Lambda\in\mathbb N$ if $\gamma_x=0$ for $|x|>\Lambda$.}
\end{definition}

\noindent\reve{The sign in the exponent is that of the Fourier symbol of a rule,
$\hat L(k)=\sum_g\ell_ge^{-\ii kg}$ of the paragraph after Lemma~\ref{lem:qca-circle}; in this chapter the coefficients
of a rule are indexed by $g\in\Z$, because the letter $m$ is needed for the mass in
Example~\ref{ex:gs-leapfrog}. \revs{Cell data need not satisfy the state condition.} The vacuum
form of Example~\ref{ex:gs-vacuum} has the cells $\gamma_0=1$ and $\gamma_x=0$ for $x\neq0$, so
that $\hat\gamma(k)=1$, and the squeezed form has $\gamma_0=\operatorname{diag}(1/\vartheta,\vartheta)$ and
$\gamma_x=0$ otherwise; both have range $0$.}

\begin{lemma}[Cell forms]\label{lem:gs-cells}
\reve{Let $(\gamma_x)$ be cell data.
\begin{enumerate}[label=(\alph*)]
\item The covariance form is a well-defined symmetric bilinear form on $V$, invariant under
translations.
\item The series for $\hat\gamma(k)$ converges absolutely, the function
$\hat\gamma$ is continuous, and $\hat\gamma(k)$ is Hermitian with
$\hat\gamma(-k)=\overline{\hat\gamma(k)}$.
\end{enumerate}}
\end{lemma}

\begin{proof}\leavevmode
\reve{\begin{enumerate}[label=(\alph*)]
\item The double sum is finite because $u$ and $w$ are finitely supported, and it is bilinear. It
is symmetric because
$\gamma(w,u)=\sum_{x,y}u(y)^{\mathsf T}\gamma_{x-y}^{\mathsf T}w(x)=\sum_{x,y}u(y)^{\mathsf
T}\gamma_{y-x}w(x)=\gamma(u,w)$, and translation invariance follows from the substitution
$x\mapsto x+a$, $y\mapsto y+a$.
\item The terms are bounded by $\|\gamma_x\|$, so that the Weierstrass test gives absolute and uniform
convergence of continuous functions, and hence continuity. Since the $\gamma_x$ are real,
$\hat\gamma(k)^\dagger=\sum_x\gamma_x^{\mathsf T}e^{\ii kx}=\sum_x\gamma_{-x}e^{\ii kx}=\hat\gamma(k)$
and $\overline{\hat\gamma(k)}=\sum_x\gamma_xe^{\ii kx}=\hat\gamma(-k)$.
\end{enumerate}}
\end{proof}

\noindent\reve{The convergence in (b) is uniform in $k$, and each term is $2\pi$-periodic, so that $\hat\gamma$ is a continuous function on the circle.}

\noindent\reve{The state condition concerns complex combinations of vectors, that is, finitely
supported functions $\mu\colon\Z\to\C^{2n}$, with Fourier transform
$\hat\mu(k)=\sum_x\mu(x)e^{-\ii kx}$. By the discussion after Definition~\ref{def:gs-statecond},
$(\gamma,\sigma)$ satisfies the state condition if and only if $Q(\mu)\ge0$ for every such $\mu$,
where
\[Q(\mu)=\sum_{x,y\in\Z}\mu(x)^\dagger a_{x-y}\,\mu(y),\qquad a_g=\gamma_g+\ii\delta_{g,0}\Omega,\]
and $\sum_ga_ge^{-\ii kg}=\hat\gamma(k)+\ii\Omega$.}

\begin{lemma}[Parseval identity]\label{lem:gs-parseval}
\reve{For cell data and every finitely supported $\mu\colon\Z\to\C^{2n}$,
\[Q(\mu)=\frac1{2\pi}\int_{-\pi}^{\pi}\hat\mu(k)^\dagger\bigl(\hat\gamma(k)+\ii\Omega\bigr)\hat\mu(k)\,dk .\]}
\end{lemma}

\begin{proof}
\reve{For fixed $k$,
\[\hat\mu(k)^\dagger\bigl(\hat\gamma(k)+\ii\Omega\bigr)\hat\mu(k)
=\sum_{x,y}\sum_g\mu(x)^\dagger a_g\mu(y)\,e^{\ii k(x-y-g)},\]
a finite sum over $x,y$ of series in $g$ that converge uniformly in $k$, the terms being bounded by
$|\mu(x)|\,|\mu(y)|\,\|a_g\|$. We may therefore integrate the series term by term, and since
$\frac1{2\pi}\int_{-\pi}^\pi e^{\ii kj}\,dk=\delta_{j,0}$ for $j\in\Z$, exactly the terms with
$g=x-y$ remain, whose sum is $Q(\mu)$.}
\end{proof}

\begin{theorem}[Fourier criterion]\label{thm:gs-fourier}
\reve{Let $(\gamma_x)$ be cell data with covariance form $\gamma$.
\begin{enumerate}[label=(\alph*)]
\item If $\hat\gamma(k)+\ii\Omega\ge0$ for every $k\in\R$, then $(\gamma,\sigma)$ satisfies the
state condition.
\item If $(\gamma_x)$ has finite range and $(\gamma,\sigma)$ satisfies the state condition, then
$\hat\gamma(k)+\ii\Omega\ge0$ for every $k\in\R$.
\end{enumerate}}
\end{theorem}

\begin{proof}\leavevmode
\reve{\begin{enumerate}[label=(\alph*)]
\item The integrand in Lemma~\ref{lem:gs-parseval} is non-negative for every $k$, hence so is
$Q(\mu)$.
\item Suppose that $v^\dagger\bigl(\hat\gamma(k_0)+\ii\Omega\bigr)v=-\varepsilon<0$ for some $k_0$
and $v\in\C^{2n}$, and let $\Lambda$ be the range. For $N>\Lambda$ let $\mu_N(x)=e^{\ii k_0x}v$ for $0\le x<N$
and $\mu_N(x)=0$ otherwise. The number of pairs $(x,y)\in\{0,\dots,N-1\}^2$ with $x-y=g$ is
$N-|g|$ for $|g|<N$, and $a_g=0$ for $|g|>\Lambda$, so that
\[Q(\mu_N)=\sum_{|g|\le\Lambda}(N-|g|)\,e^{-\ii k_0g}\,v^\dagger a_gv
=-N\varepsilon-\sum_{|g|\le\Lambda}|g|\,e^{-\ii k_0g}\,v^\dagger a_gv .\]
The last sum is bounded in absolute value by $C=\sum_{|g|\le\Lambda}|g|\,|v^\dagger a_gv|$, which does
not depend on $N$, and for $N>\max(\Lambda,C/\varepsilon)$ we obtain $Q(\mu_N)\le-N\varepsilon+C<0$,
contradicting the state condition.
\end{enumerate}}
\end{proof}

\begin{remark}[Necessity for all cell data]\label{rem:gs-fourier-general}
\reve{The hypothesis of finite range in part~(b) can be dropped: for all cell data, the state
condition holds if and only if $\hat\gamma(k)+\ii\Omega\ge0$ for every $k\in\R$. Only necessity
needs proof. Assume the state condition, fix $k_0\in\R$ and $v\in\C^{2n}$, and define $\mu_N$ as in
the proof of~(b) for every $N\ge1$. The state condition gives $Q(\mu_N)\ge0$ for every $N$, and
counting the pairs $(x,y)$ with $x-y=g$ as there, now without any restriction on $g$,
\[\frac{Q(\mu_N)}N=\sum_{|g|<N}\Bigl(1-\frac{|g|}N\Bigr)e^{-\ii k_0g}\,v^\dagger a_gv .\]
The $g$-th term tends to $e^{-\ii k_0g}\,v^\dagger a_gv$ as $N\to\infty$ and is bounded in absolute
value by $|v^\dagger a_gv|\le|v|^2\|a_g\|$, which is summable over $g$ because the cells are
absolutely summable. By dominated convergence for series, $Q(\mu_N)/N$ tends to
$\sum_ge^{-\ii k_0g}\,v^\dagger a_gv=v^\dagger\bigl(\hat\gamma(k_0)+\ii\Omega\bigr)v$, which is
therefore a limit of non-negative numbers and non-negative. As $v$ and $k_0$ were arbitrary,
$\hat\gamma(k)+\ii\Omega\ge0$ for every $k$, and together with part~(a) the criterion is an
equivalence for all cell data.}
\end{remark}

\begin{remark}[The other sign]\label{rem:gs-minus}
\reve{By Lemma~\ref{lem:gs-cells}(b), $\hat\gamma(k)-\ii\Omega=\overline{\hat\gamma(-k)+\ii\Omega}$,
and the complex conjugate of a positive semidefinite matrix is positive semidefinite. If
$\hat\gamma(k)+\ii\Omega\ge0$ for every $k$, therefore, both $\hat\gamma(k)\pm\ii\Omega$ are positive
semidefinite at every $k$, and so is their average $\hat\gamma(k)$.}
\end{remark}

\subsection{Saturation of the state condition}

\noindent\reve{Purity of a state is not defined in these notes, and nothing in this subsection
proves that a state is pure. We prove the equivalence of three matrix conditions on the
symbol which, by analogy with finitely many modes, are candidates for the characterization of the
pure translation-invariant Gaussian states.}

\noindent\motiv{\reve{A Gaussian state of finitely many modes is pure exactly when its covariance
matrix saturates the uncertainty relation, which for one mode means $\det\gamma_B=1$. For a
translation-invariant state on the chain with one mode per site, the Fourier criterion suggests
that saturation should be read at every wave number, where it can be expressed in three ways,
namely that $\hat\gamma(k)+\ii\Omega$ is singular, that $\det\hat\gamma(k)=1$, and that $\Omega\hat
\gamma(k)$ squares to $-1$, as $\Omega\gamma_B$ does for a pure state of finitely many modes. For a single $2\times2$
matrix the first is weaker than the others unless both signs of $\ii\Omega$ are taken into
account.}}

\begin{lemma}[Two-by-two matrices]\label{lem:gs-twobytwo}
\reve{Let $n=1$ and let $A=\begin{psmallmatrix}a&c\\\bar c&b\end{psmallmatrix}$ be a Hermitian
$2\times2$ matrix with $A+\ii\Omega\ge0$ and $A-\ii\Omega\ge0$. Then $\det A\ge1+2|\operatorname{Im}
c|$, and the following are equivalent:
\begin{enumerate}[label=(\alph*)]
\item $(\Omega A)^2=-1$;
\item $\det A=1$;
\item $A+\ii\Omega$ and $A-\ii\Omega$ are both singular.
\end{enumerate}}
\end{lemma}

\begin{proof}
\reve{Since
$A\pm\ii\Omega=\begin{psmallmatrix}a&c\pm\ii\\\bar c\mp\ii&b\end{psmallmatrix}$ and
$(c\pm\ii)(\bar c\mp\ii)=|c|^2+1\pm2\operatorname{Im}c$,
\begin{equation}\label{eq:gs-detpm}
\det(A\pm\ii\Omega)=\det A-1\mp2\operatorname{Im}c ,
\end{equation}
and both determinants are non-negative by hypothesis, which gives the bound.}

\reve{By \eqref{eq:gs-detpm}, condition (c) says $\det A=1+2\operatorname{Im}c=1-2\operatorname{Im}c$,
that is, $\operatorname{Im}c=0$ and $\det A=1$, so that (c) implies (b); conversely, $\det A=1$
forces $\operatorname{Im}c=0$ by the bound, and then both determinants vanish, which proves
(b)$\Leftrightarrow$(c).}

\reve{For (a)$\Leftrightarrow$(b), put $M=\Omega A=\begin{psmallmatrix}\bar
c&b\\-a&-c\end{psmallmatrix}$, so that $\operatorname{tr}M=-2\ii\operatorname{Im}c$ and $\det
M=\det A$; by the Cayley--Hamilton theorem, $M^2=(\operatorname{tr}M)M-(\det M)1$. If (b) holds,
then $\operatorname{Im}c=0$, as just shown, so that $\operatorname{tr}M=0$ and $M^2=-1$. Suppose
conversely that $M^2=-1$. If $\operatorname{tr}M\neq0$, the relation $(\operatorname{tr}M)M=(\det
M-1)1$ makes $M$ a scalar $\lambda1$ with $\lambda^2=-1$, and $A=\Omega^{-1}M=-\lambda\Omega$; for
$\lambda=\ii$ this gives $A-\ii\Omega=-2\ii\Omega$, and for $\lambda=-\ii$ it gives
$A+\ii\Omega=2\ii\Omega$, and either matrix has the eigenvalues $2$ and $-2$, contrary to the
hypothesis. Hence $\operatorname{tr}M=0$, and $M^2=-(\det M)1=-1$ gives $\det A=1$.}
\end{proof}

\noindent\reve{Both signs are needed for (b)$\Rightarrow$(a): the matrix
$A=\begin{psmallmatrix}\sqrt2&-\ii\\\ii&\sqrt2\end{psmallmatrix}$ has $\det A=1$ and
$A+\ii\Omega=\sqrt2\cdot1\ge0$, but $(\Omega A)^2=\begin{psmallmatrix}-3&2\sqrt2\,\ii\\
-2\sqrt2\,\ii&-3\end{psmallmatrix}$. The example does not contradict the lemma, because it
violates its hypothesis: $A-\ii\Omega=\begin{psmallmatrix}\sqrt2&-2\ii\\2\ii&\sqrt2\end{psmallmatrix}$
has determinant $-2$ and is therefore not positive semidefinite. For a real matrix $A$ the two hypotheses coincide, and for the
symbol of cell data they coincide by Remark~\ref{rem:gs-minus}.}

\begin{corollary}[Saturation at every wave number]\label{cor:gs-purity}
\reve{Let $n=1$ and let $(\gamma_x)$ be cell data with $\hat\gamma(k)+\ii\Omega\ge0$ for
every $k$. The following are equivalent:
\begin{enumerate}[label=(\alph*)]
\item $\hat\gamma(k)+\ii\Omega$ is singular for every $k$;
\item $\det\hat\gamma(k)=1$ for every $k$;
\item $(\Omega\hat\gamma(k))^2=-1$ for every $k$.
\end{enumerate}}
\end{corollary}

\begin{proof}
\reve{By Remark~\ref{rem:gs-minus}, Lemma~\ref{lem:gs-twobytwo} applies to $A=\hat\gamma(k)$ for every
$k$, and it gives (b)$\Leftrightarrow$(c) at each $k$ as well as (b)$\Rightarrow$(a). For
(a)$\Rightarrow$(b), note that $\hat\gamma(-k)+\ii\Omega=\overline{\hat\gamma(k)-\ii\Omega}$, so
that (a) at $-k$ says that $\hat\gamma(k)-\ii\Omega$ is singular; both $\hat\gamma(k)\pm\ii\Omega$
are then singular, and Lemma~\ref{lem:gs-twobytwo} gives $\det\hat\gamma(k)=1$.}
\end{proof}

\noindent\reve{The equivalence holds for the conditions over all $k$ together and not wave number by
wave number: at a single $k$, singularity of $\hat\gamma(k)+\ii\Omega$ does not imply
$\det\hat\gamma(k)=1$ when $\operatorname{Im}\hat\gamma(k)_{12}>0$, which can happen when the cells
$\gamma_x$ with $x\neq0$ are not symmetric matrices. For example, the cell data
$\gamma_0=2\cdot1$, $\gamma_1=\tfrac12\begin{psmallmatrix}0&-1\\1&0\end{psmallmatrix}$,
$\gamma_{-1}=\gamma_1^{\mathsf T}$ and $\gamma_x=0$ otherwise have the symbol
$\hat\gamma(k)=\begin{psmallmatrix}2&\ii\sin k\\-\ii\sin k&2\end{psmallmatrix}$, so that
$\det(\hat\gamma(k)+\ii\Omega)=4-(1+\sin k)^2\ge0$ and, the diagonal being positive,
$\hat\gamma(k)+\ii\Omega\ge0$ for every $k$; at $k=\pi/2$ this matrix is singular, whereas
$\det\hat\gamma(\pi/2)=3$.}

\noindent\reve{Both product examples satisfy the conditions of Corollary~\ref{cor:gs-purity}. For the vacuum, $1+\ii\Omega$ has the
eigenvalues $0$ and $2$ and $(\Omega\cdot1)^2=\Omega^2=-1$; the squeezed form has the constant symbol
$\operatorname{diag}(1/\vartheta,\vartheta)$, whose determinant is $1$.}

\subsection{Evolution}

\noindent\motiv{\reve{A symplectic rule acts on states by pullback along its automorphism, and on
the data of a Gaussian state by composition with the rule, so that Gaussian states evolve into
Gaussian states. On a finite region an on-site rule acts on the covariance matrix by the symplectic
congruence familiar from finitely many modes, and on the chain the evolution of the cell data is
a congruence of the symbol at each wave number. Until the symbol is computed, $G$ is again an arbitrary abelian
group.}}

\begin{proposition}[Evolved states]\label{prop:gs-evolve}
\reve{Let $L$ be a symplectic rule and $\omega$ a state.
\begin{enumerate}[label=(\alph*)]
\item $\omega\circ\alpha_L$ is a state with $\chi_{\omega\circ\alpha_L}=\chi_\omega\circ L$.
\item If $(\omega,\gamma,d)$ is a Gaussian state, then so is
$\bigl(\omega\circ\alpha_L,\ \gamma(L\cdot,L\cdot),\ d\circ L\bigr)$.
\item If $\omega$ is translation invariant, so is $\omega\circ\alpha_L$.
\end{enumerate}}
\end{proposition}

\begin{proof}\leavevmode
\reve{\begin{enumerate}[label=(\alph*)]
\item Proposition~\ref{prop:gs-pullback}(b) with $\alpha=\alpha_L$ and $Q=L$.
\item Since $L$ is linear, $\gamma(L\cdot,L\cdot)$ is a symmetric bilinear form and $d\circ L$ a
linear form, and by (a)
$\chi_{\omega\circ\alpha_L}(u)=\exp\bigl(\ii d(Lu)-\tfrac14\gamma(Lu,Lu)\bigr)$.
\item By Proposition~\ref{prop:alpha}(c), $\alpha_L\alpha_a=\alpha_a\alpha_L$, so that
$(\omega\circ\alpha_L)\circ\alpha_a=(\omega\circ\alpha_a)\circ\alpha_L=\omega\circ\alpha_L$.
\end{enumerate}}
\end{proof}

\noindent\motiv{\reve{The automorphisms of rules compose as the rules do, so that iterating an
automaton on a state amounts to multiplying symbols.}}

\begin{proposition}[Composition of rules]\label{prop:gs-hom}
\reve{The symplectic symbols form a submonoid of the matrices over the group algebra $\R[G]$, and
$\hat L\mapsto\alpha_L$ is a monoid homomorphism from it into the $*$-automorphisms of
$\CCR(V,\sigma)$ under composition, which restricts to a group homomorphism on the units.}
\end{proposition}

\begin{proof}
\reve{The identity is symplectic, and if $\hat L^{\star\mathsf T}\Omega\hat L=\Omega$ and $\hat
N^{\star\mathsf T}\Omega\hat N=\Omega$, then, since $\hat M\mapsto\hat M^{\star\mathsf T}$ reverses
products (the involution is a ring homomorphism of the commutative ring $\R[G]$, and
transposition reverses products), $(\hat L\hat N)^{\star\mathsf T}\Omega\hat L\hat N=\hat
N^{\star\mathsf T}(\hat L^{\star\mathsf T}\Omega\hat L)\hat N=\Omega$. The homomorphism property is
Proposition~\ref{prop:alpha}(b), and a monoid homomorphism maps units to units.}
\end{proof}

\noindent\reve{By Corollary~\ref{cor:qca-group}, every symplectic symbol is invertible, with a
symplectic inverse, so that the submonoid of symplectic symbols is a group and the statement about
units applies to all of it.}

\noindent\motiv{\reve{On a finite region, the simplest rules are those acting at each site by the
same symplectic matrix, and for them the evolution of the covariance matrix can be written down
explicitly.}}

\begin{corollary}[On-site evolution]\label{cor:gs-onsite}
\reve{Let $S$ be a real $2n\times2n$ matrix with $S^{\mathsf T}\Omega S=\Omega$, let $L_S$ be the
on-site rule $(L_Su)(x)=Su(x)$, and let $O$ be a finite region. For every bilinear form $\gamma$,
\[\cov_O\bigl(\gamma(L_S\cdot,L_S\cdot)\bigr)=T_O\,\cov_O(\gamma)\,T_O^{\mathsf T},\]
where $T_O$ is the matrix of $u\mapsto S^{-1}u$ applied at each site of $O$, written in the
coordinates of Definition~\ref{def:gs-covmat}. If $\gamma$ is symmetric and satisfies the state
condition on $O$, so does $\gamma(L_S\cdot,L_S\cdot)$.}
\end{corollary}

\begin{proof}
\reve{Let $S_O$ be the matrix of $L_S$ on $\loc(O)$ in the coordinates of
Lemma~\ref{lem:gs-reorder}(b), so that $R_O(\gamma(L_S\cdot,L_S\cdot))=S_O^{\mathsf T}R_O(\gamma)S_O$
and $S_O^{\mathsf T}\Omega_{n_O}S_O=\Omega_{n_O}$, the latter because $L_S$ preserves $\sigma$ and
$R_O(\sigma)=\Omega_{n_O}$. From $S^{\mathsf T}\Omega S=\Omega$ and $\Omega\Omega^{\mathsf T}=1$ it
follows that $\Omega^{\mathsf T}S^{\mathsf T}\Omega=S^{-1}$, and site by site the same computation
gives $\Omega_{n_O}^{\mathsf T}S_O^{\mathsf T}\Omega_{n_O}=S_O^{-1}=T_O$. Inserting
$\Omega_{n_O}\Omega_{n_O}^{\mathsf T}=1$ twice,
\[\Omega_{n_O}^{\mathsf T}S_O^{\mathsf T}R_O(\gamma)S_O\Omega_{n_O}
=\bigl(\Omega_{n_O}^{\mathsf T}S_O^{\mathsf T}\Omega_{n_O}\bigr)
\bigl(\Omega_{n_O}^{\mathsf T}R_O(\gamma)\Omega_{n_O}\bigr)
\bigl(\Omega_{n_O}^{\mathsf T}S_O\Omega_{n_O}\bigr)=T_O\,\cov_O(\gamma)\,T_O^{\mathsf T},\]
which is the formula. For the last statement, $L_S$ maps $\loc(O)$ into itself and preserves
$\sigma$, so that for $u,w\in\loc(O)$ the real-pair form for $\gamma$ at $L_Su,L_Sw$ gives
$\gamma(L_Su,L_Su)\ge0$ and $\sigma(u,w)^2=\sigma(L_Su,L_Sw)^2\le\gamma(L_Su,L_Su)\gamma(L_Sw,L_Sw)$.}
\end{proof}

\noindent\reve{The covariance matrix transforms with $S^{-1}$ and not with $S$. The state
$\omega\circ\alpha_{L_S}$ is the evolution in the Heisenberg picture, in which the Weyl operators are
moved by $S$, and the moments of the canonical coordinates are then moved by the inverse. The
two congruences agree only in special cases, for instance for a rotation acting on the vacuum, where
both leave the identity unchanged, and in general they differ: for one site and one mode,
$S=\operatorname{diag}(2,1/2)$ takes the vacuum covariance $1$ to $S^{-1}S^{-\mathsf
T}=\operatorname{diag}(1/4,4)$, whereas $SS^{\mathsf T}=\operatorname{diag}(4,1/4)$; and a rotation
$S$ by an angle $\varphi$ with $\sin2\varphi\neq0$, acting on the squeezed covariance
$\operatorname{diag}(\vartheta,1/\vartheta)$ with $\vartheta\neq1$, gives two matrices whose
off-diagonal entries have opposite signs. By Proposition~\ref{prop:gs-admissible}, the admissibility of
$T_O\cov_O(\gamma)T_O^{\mathsf T}$ is the familiar invariance of admissibility under symplectic
congruences.}

\begin{example}[On-site rotations]\label{ex:gs-rotation}
\reve{Let $n=1$ and let $R_\varphi$ be the rotation of Example~\ref{ex:ccr-onsite}; the letter
$R$ here denotes a rotation matrix and not a restricted matrix $R_O$. Since
$R_\varphi^{\mathsf T}R_\varphi=1$, the vacuum form is invariant under the on-site rule,
$\gamma(L_{R_\varphi}u,L_{R_\varphi}w)=\gamma(u,w)$. By Proposition~\ref{prop:gs-evolve}(b), the free
evolution of uncoupled oscillators therefore takes a Gaussian state with the vacuum covariance to one
with the same covariance and a rotated displacement; on a finite region this is
Corollary~\ref{cor:gs-onsite} with $T_O\cdot1\cdot T_O^{\mathsf T}=1$.}
\end{example}

\noindent\motiv{\reve{On the chain, the covariance form of translation-invariant data evolves into
another such form, and its cells and symbol can be computed from those of the rule.}}

\begin{theorem}[Evolution of the covariance]\label{thm:gs-evolution}
\reve{Let $G=\Z$, let $L$ be a rule with coefficients $\ell_g$, finitely many nonzero, and let
$(\gamma_x)$ be cell data with covariance form $\gamma$. Put
\[\gamma'_j=\sum_{g,g'}\ell_g^{\mathsf T}\,\gamma_{j+g-g'}\,\ell_{g'}\qquad(j\in\Z).\]
Then $(\gamma'_j)$ is cell data whose covariance form is $\gamma(L\cdot,L\cdot)$, and
\[\hat\gamma'(k)=\hat L(-k)^{\mathsf T}\,\hat\gamma(k)\,\hat L(k)=\hat L(k)^\dagger\,\hat\gamma(k)\,\hat L(k).\]
If $L$ is reflection symmetric, $\ell_{-g}=\ell_g$ for all $g$, then $\hat L(-k)=\hat L(k)$ and
$\hat\gamma'(k)=\hat L(k)^{\mathsf T}\hat\gamma(k)\hat L(k)$.}
\end{theorem}

\begin{proof}
\reve{The cells $\gamma'_j$ are finite combinations of shifts of $(\gamma_x)$ and are therefore
absolutely summable. They are symmetric, since $(\gamma'_j)^{\mathsf T}=\sum_{g,g'}\ell_{g'}^{\mathsf
T}\gamma_{-j-g+g'}\ell_g$, which after exchanging the names $g$ and $g'$ is $\gamma'_{-j}$. For the
form,
\[\gamma(Lu,Lw)=\sum_{x,y}\sum_{g,g'}u(x-g)^{\mathsf T}\ell_g^{\mathsf T}\gamma_{x-y}\ell_{g'}w(y-g'),\]
and the substitution $x'=x-g$, $y'=y-g'$ turns $\gamma_{x-y}$ into $\gamma_{x'-y'+g-g'}$, so that the
inner sum over $g,g'$ is $\gamma'_{x'-y'}$. For the symbol, the substitution $p=j+g-g'$ gives
$e^{-\ii kj}=e^{\ii kg}e^{-\ii kp}e^{-\ii kg'}$, and since the sums over $g,g'$ are finite and the
sum over $p$ converges absolutely, we may rearrange the series into
$\bigl(\sum_g\ell_g^{\mathsf T}e^{\ii kg}\bigr)\hat\gamma(k)\bigl(\sum_{g'}\ell_{g'}e^{-\ii kg'}\bigr)
=\hat L(-k)^{\mathsf T}\hat\gamma(k)\hat L(k)$. As the $\ell_g$ are real,
$\hat L(-k)^{\mathsf T}=\overline{\hat L(k)}^{\mathsf T}=\hat L(k)^\dagger$, and the last statement
is immediate.}
\end{proof}

\noindent\reve{The second form shows that the evolution is a congruence at each wave number. If $L$ is
symplectic, then $\hat L(k)^\dagger\Omega\hat L(k)=\Omega$ by Example~\ref{ex:qca-modes}, so that
$\hat\gamma'(k)+\ii\Omega=\hat L(k)^\dagger\bigl(\hat\gamma(k)+\ii\Omega\bigr)\hat L(k)$, and the
evolution preserves the Fourier criterion, in accordance with Proposition~\ref{prop:gs-evolve}(b)
and Theorem~\ref{thm:gs-necessity}. The formula with $\hat L(-k)^{\mathsf T}=\hat L(k)^\dagger$ is
the general form of the evolution law, and without reflection symmetry the transpose alone cannot
replace the conjugate transpose: for the rule with $\ell_0=1$, $\ell_1=\begin{psmallmatrix}0&1\\0&0
\end{psmallmatrix}$ and all other coefficients zero, the vacuum symbol is taken to $\hat L(k)^\dagger
\hat L(k)=\begin{psmallmatrix}1&e^{-\ii k}\\e^{\ii k}&2\end{psmallmatrix}$, whereas $\hat
L(k)^{\mathsf T}\hat L(k)=\begin{psmallmatrix}1&e^{-\ii k}\\e^{-\ii k}&1+e^{-2\ii k}\end{psmallmatrix}$
is not even Hermitian.}

\begin{example}[The leapfrog applied to the vacuum]\label{ex:gs-leapfrog}
\reve{The leapfrog $F$ of Example~\ref{ex:qca-leapfrog-rule} has the coefficients
$\ell_0=\begin{psmallmatrix}0&1\\-1&-m^2\end{psmallmatrix}$ and
$\ell_{\pm1}=\begin{psmallmatrix}0&0\\0&1\end{psmallmatrix}$, so that it is reflection symmetric,
with $\hat F(k)=\begin{psmallmatrix}0&1\\-1&c_k\end{psmallmatrix}$ and $c_k=2\cos k-m^2$, where $m$
is the mass. By Proposition~\ref{prop:time-direction}, the rule that moves Cauchy data forward is $F^{-1}$, and
reading $\alpha_{F^{-1}}$ as the forward time step of the lattice field is the heuristic reading that
follows that proposition. The rule $F^{-1}$ has the symbol $\hat
F^{-1}(k)=\begin{psmallmatrix}c_k&-1\\1&0\end{psmallmatrix}$, that is, the coefficients
$\ell_0=\begin{psmallmatrix}-m^2&-1\\1&0\end{psmallmatrix}$ and
$\ell_{\pm1}=\begin{psmallmatrix}1&0\\0&0\end{psmallmatrix}$, and it is again reflection symmetric.
By Theorem~\ref{thm:gs-evolution}, one forward step takes the vacuum symbol $\hat\gamma(k)=1$ to
\[\hat\gamma'(k)=\hat F^{-1}(k)^{\mathsf T}\hat F^{-1}(k)=\begin{pmatrix}1+c_k^2&-c_k\\-c_k&1\end{pmatrix},\]
whose cells are
\[\gamma'_0=\begin{pmatrix}3+m^4&m^2\\m^2&1\end{pmatrix},\qquad
\gamma'_{\pm1}=\begin{pmatrix}-2m^2&-1\\-1&0\end{pmatrix},\qquad
\gamma'_{\pm2}=\begin{pmatrix}1&0\\0&0\end{pmatrix},\]
and zero otherwise, as the formula $\gamma'_j=\sum_g\ell_g^{\mathsf T}\ell_{g+j}$ for the
coefficients of $F^{-1}$ confirms directly together with $1+c_k^2=3+m^4-4m^2\cos k+2\cos2k$. Since
$\det\hat\gamma'(k)=1$ for every $k$, the evolved covariance again satisfies the equivalent conditions of
Corollary~\ref{cor:gs-purity}, while its correlations, absent in the vacuum, now reach to the second
neighbour: a single step of range~$1$ creates correlations of range~$2$, because the covariance is
quadratic in the rule. The backward step $\alpha_F$ gives in the same way
$\hat F(k)^{\mathsf T}\hat F(k)=\begin{psmallmatrix}1&-c_k\\-c_k&1+c_k^2\end{psmallmatrix}$, the same
matrix with its two diagonal entries exchanged.}
\end{example}

\subsection{Notes and references}\label{sec:gs-notes}

\begin{enumerate}
\item \reve{\emph{States and Gaussian states.} In~\cite[p.~77]{krueger}, states are positive,
normalized linear functionals, equivalently described by their characteristic functions, and a
state is Gaussian when its characteristic function is
$\exp\bigl(-\tfrac14\gamma(\xi,\xi)+\ii\sum_x\xi_x^{\mathsf T}d_x\bigr)$, with the same factor
$\tfrac14$ as Definition~\ref{def:gs-gaussian}. The thesis does not discuss the existence of
Gaussian states on the lattice CCR algebra; Remarks~\ref{rem:gs-existence}
and~\ref{rem:gs-converse} record what is involved.}
\item \reve{\emph{State condition.} Condition (a) of Definition~\ref{def:gs-statecond} is the
thesis's (4.5), $\gamma(\bar\mu,\mu)+\ii\sigma(\bar\mu,\mu)\ge0$, and $Q$ in
Subsection~\ref{sec:gs-fourier} is its form (4.6)~\cite[pp.~77--78]{krueger}. The thesis obtains the
necessity of the state condition ``by a direct generalization of the argument leading to (2.22)'';
the proof of Theorem~\ref{thm:gs-necessity} is a different, derivative-free argument. The thesis
works throughout with $\gamma+\ii\sigma$, which corresponds to the finite-mode condition written
with $\gamma_B-\ii\Omega$; by Proposition~\ref{prop:gs-statecond} and
Remark~\ref{rem:gs-normalization} the sign is immaterial.}
\item \reve{\emph{Finite regions.} The thesis's correlation matrix of a region is the matrix of
the form $\gamma$ on the site vectors, ordered site by site, whereas the finite-mode covariance
matrix is ordered positions first and is conjugated by $\Omega$; the normalization, fixed by the
factor $\tfrac14$, is the same.
Definition~\ref{def:gs-covmat} and Proposition~\ref{prop:gs-admissible} make the translation between
the two explicit; they have no counterpart in the thesis.}
\item \reve{\emph{Fourier criterion.} The Fourier convention is the thesis's (4.14), the cell indexing
$\gamma_{x-y}$ is that of (4.4), and Lemma~\ref{lem:gs-parseval} is (4.16)~\cite[pp.~77,
81]{krueger}. Theorem~\ref{thm:gs-fourier} is Lemma~4.5(i) of the thesis~\cite[p.~82]{krueger},
which asserts both directions for summable cell data and argues the direction (b) with the
``sharply peaked Fourier transform of a flat Gaussian''~\cite[p.~81]{krueger}; here we prove (b)
by the same idea with a rectangular window, first for finite range and then, in
Remark~\ref{rem:gs-fourier-general}, for all summable cell data.}
\item \reve{\emph{Purity.} Lemma~4.5(ii) of~\cite[p.~82]{krueger} states, for one mode per site, that
the correlation function ``corresponds to a pure Gaussian state if and only if
$\hat\gamma(k)+\ii\sigma_s$ is not strictly positive for any $k\in[-\pi,\pi]$, i.e.\ if
$(\sigma_s\hat\gamma(k))^2=-1$''. Corollary~\ref{cor:gs-purity} proves the equivalence of the two
matrix conditions, which holds for all $k$ together; the ``i.e.''\ is not an equivalence of the
pointwise conditions, as the remark after the corollary explains. The first clause, the
characterization of purity, is not proved here. The vacuum is the thesis's ``coherent product
state'' $\hat\gamma_0(k)=1$.}
\item \reve{\emph{Evolution.} The cells $\gamma'_j$ of Theorem~\ref{thm:gs-evolution} are the thesis's
(4.17a), $\gamma_t(x)=\sum_{y,z}\Gamma_t(y)^{\mathsf T}\gamma_0(x+y-z)\Gamma_t(z)$, and the printed
symbol formula (4.17b), \reve{$\hat\gamma_t(k)=\overline{\hat\Gamma_t^{\mathsf T}(k)}\,\hat\gamma_0(k)\,\hat\Gamma_t(k)$
\cite[p.~82]{krueger}, with the conjugate transpose, is the general form $\hat\gamma_t(k)=\hat\Gamma_t(k)^\dagger\hat\gamma_0(k)\hat\Gamma_t(k)$ of Theorem~\ref{thm:gs-evolution}. For the rules of the thesis, which are reflection symmetric by (4.12)~\cite[p.~80]{krueger} and have real symbols, the conjugate transpose is the transpose.}}
\item \reve{\emph{Existence of states.} Quasi-free states of the CCR algebra, of which the Gaussian
states of this chapter are the lattice instance, were introduced by Manuceau and
Verbeure~\cite{manuceau-verbeure}. The theorem that a normalized function with positive
semidefinite twisted kernels is the characteristic function of a unique state is due to Manuceau,
Sirugue, Testard and Verbeure~\cite{mstv}; see also~\cite[\S5.2]{bratteli-robinson}. The quantum
Bochner theorem of Remark~\ref{rem:gs-bochner} is due to Kastler~\cite{kastler} and to Loupias and
Miracle-Sole~\cite{lms}.}
\end{enumerate}
